\documentclass[12pt]{article}
\usepackage[margin=0.65in]{geometry}

\newcommand{\duedate}{}
\newcommand{\assignment}{End-To-End Memory Networks}

\newcommand{\name}{Aksel Kretsinger-Walters}
\newcommand{\email}{}

\makeatletter
\def\input@path{{../}{../../}{../../../}}
\makeatother
\input{pset_template_CLMM.tex}

\begin{document}
\psetheader

\section*{End-To-End Memory Networks}

\paragraph{Motivation.}

Many reasoning tasks require models to consult stored knowledge before
producing an answer. Traditional neural networks compress all information
into hidden states, which limits their ability to retrieve specific facts.

End-to-End Memory Networks introduce an explicit memory component
that can be queried through attention.

\paragraph{Architecture.}

The model consists of three main components:

\begin{itemize}
		\item An input representation
		\item A memory module
		\item An output module
\end{itemize}

The memory stores a set of vectors representing facts or observations:

\[
		m_i \in \mathbb{R}^d
\]

A query vector $u$ represents the current question.

\paragraph{Memory Attention.}

The model computes attention weights over memory entries:

\[
		p_i = \text{Softmax}(u^\top m_i)
\]

These weights determine which facts are relevant.

A response vector is computed as

\[
		o = \sum_i p_i c_i
\]

where $c_i$ are output embeddings associated with each memory entry.

\paragraph{Multi-Hop Reasoning.}

The model can perform multiple reasoning steps (``hops'').

After each hop, the query is updated:

\[
		u_{k+1} = u_k + o_k
\]

This allows the model to iteratively refine its reasoning and combine
information from multiple memory entries.

\paragraph{Prediction.}

The final prediction is produced via

\[
		\hat{a} = \text{Softmax}(W(u_K + o_K))
\]

where $K$ is the number of hops.

\paragraph{Applications.}

The model was evaluated on the bAbI dataset, which contains
synthetic reasoning tasks such as:

\begin{itemize}
		\item question answering
		\item path finding
		\item logical deduction
\end{itemize}

Multi-hop memory access allows the model to combine several supporting facts.

\paragraph{Advantages.}

Compared to earlier memory networks, this architecture:

\begin{itemize}
		\item is fully differentiable
		\item can be trained end-to-end with gradient descent
		\item does not require supervision of memory retrieval steps
\end{itemize}

\paragraph{Limitations.}

Several challenges remain:

\begin{itemize}
		\item Memory size must remain relatively small
		\item Attention mechanisms scale poorly with large knowledge bases
		\item Performance degrades on more complex reasoning tasks
\end{itemize}

These limitations motivated later architectures such as
transformers and retrieval-augmented models.

\paragraph{Takeaway.}

End-to-End Memory Networks introduced differentiable attention over
explicit memory, enabling neural networks to retrieve and combine
facts dynamically. This work strongly influenced later attention-based
architectures and can be viewed as an early precursor to modern
transformer models.

\end{document}
